% Section 7: Discussion and Limitations (about 0.5 page). Plan: PAPER_PLAN 4 / §7.
% One model serves both tasks; shared semantic context lifts category, private spatial path keeps region competitive
% (matches at small region counts, outperforms at large). Limitations: single seed over five folds (provisional);
% small-state region is non-inferior within two points (a match, not a beat); transductive representation
% (planned training-only variant); the mobility-aware service is motivation, not evaluated.
% 2026-07-08 (P7/C4-C6): the usage sketch is now quantified (shortlist enrichment from Tbl 3 + shortlist
% compactness + near-miss distance vs the random-pair floor, all from analysis/{near_miss,shortlist_compactness}_
% RESULTS.md; four datasets, single-seed tag inline) and the closing paragraph carries the moura2025 succession
% clause. Still motivation-only: bare prediction properties, no service metric.
\section{Discussion and Limitations}
\label{sec:discussion}

The clearest reading of our results is a design one: one model is enough. A shared trunk carries the semantic context that lifts the next-category task sharply over a dedicated model at every dataset, and a private spatial path keeps the next-region task competitive: a non-inferior match (TOST, $\pm2$~pp) at the two small U.S. states, and a tested gain at the other four datasets. Across the U.S. states the region gain rises with the number of regions (Figure~\ref{fig:deltas}). Overall, the joint model matches or outperforms the dedicated single-task model on both tasks at once (Table~\ref{tab:results}): for a mobility-aware service, one model, one forward pass, anticipates both what a user will do next and where, with no task trade-off.

As a usage sketch, a service could treat the model's top-ranked next regions as an anticipatory set: at California, a shortlist of ten regions out of 8{,}501 contains the true next region 65.69 percent of the time, more than five hundred times better than
picking ten at random; at Alabama, ten out of 1{,}109 capture it 69.80 percent
of the time.  On the four datasets we measured (Alabama, Arizona, Florida, and Istanbul; a single seed over five folds), the ten shortlisted regions sit, on average, about 3 to 8 kilometers from the shortlist's own center (median across visits, from region centroids). When the top guess is wrong and the true region was seen in training, the predicted region's centroid sits a similar median of about 3 to 8 kilometers from the true region's, while two randomly drawn regions from the same map sit a median
of 20 to 241 kilometers apart. This remains motivation, not a measured service result: we report
properties of the model's own predictions, not a system's cost or a service's
outcome.

Three limits qualify these results. First, the joint numbers at California and Texas come from a single seed over five folds, so those two cells are provisional; the other four datasets average four seeds on both arms. Second, our representation is trained once over all places; visits to places
never seen during training are the single effect we cannot fully isolate. A planned follow-up trains it on each fold's training places only and embeds unseen visits directly. Third, although a mobility-aware service motivates this work, we do not build or evaluate one; a service evaluation remains the natural next study.

Finally, our representation is a fixed per-visit vector that any downstream model can consume, whether in other sequence models or in mobility tasks beyond next-category and next-region. Structural analysis of check-in networks has itself named machine
learning as a next step~\cite{moura2025mobilityaware}; learning over check-in
structure, as this study does, moves in that direction.
